% titel: Blatt 7, A3: Regel von de l'Hopital
% skript: Kap. 4 (Grenzwerte und Differentialrechnung)

% Teil I: Werkzeuge, kein Lösungsweg.

Die Regel von de l'Hôpital vereinfacht Grenzwerte von Brüchen, bei denen Zähler
und Nenner gleichzeitig gegen null oder gegen unendlich streben
(unbestimmte Formen $0/0$ bzw. $\infty/\infty$). Sie ersetzt den Quotienten
durch den Quotienten der Ableitungen.

\begin{werkzeugkasten}[Regel von de l'Hopital]
Sind $f$ und $g$ differenzierbar, $g'\neq 0$ in einer Umgebung von $x_0$
(außer evtl.\ $x_0$ selbst), und gilt
\begin{equation}
  \lim_{x\to x_0}f(x)=0,\quad \lim_{x\to x_0}g(x)=0
  \quad\text{oder}\quad
  \lim_{x\to x_0}|f(x)|=\lim_{x\to x_0}|g(x)|=\infty,
\end{equation}
so ist, falls der rechte Grenzwert existiert,
\begin{equation}
  \lim_{x\to x_0}\frac{f(x)}{g(x)}
  =\lim_{x\to x_0}\frac{f'(x)}{g'(x)}.
\end{equation}
Die Regel lässt sich mehrfach anwenden und gilt auch für einseitige
Grenzwerte und $x_0=\pm\infty$. Sie ist nur für unbestimmte Formen gültig.

\textit{Skript:} Kap. 4, Grenzwerte von Funktionen; Herkunft über den
verallgemeinerten Mittelwertsatz.
\end{werkzeugkasten}

\begin{werkzeugkasten}[Typische unbestimmte Formen und Umformungen]
Nicht immer liegt ein Bruch vor. Weitere Formen sind $0\cdot\infty$,
$\infty-\infty$, $1^\infty$, $0^0$, $\infty^0$. Sie werden durch Umformen
(kehrwert bilden, Logarithmus) auf $0/0$ oder $\infty/\infty$ gebracht, bevor
de l'Hôpital anwendbar wird. Vor dem Ableiten die Voraussetzung prüfen.

\textit{Skript:} Kap. 4.
\end{werkzeugkasten}

\begin{achtung}[Wichtig]
\begin{itemize}
  \item Erst feststellen, ob wirklich $0/0$ oder $\infty/\infty$ vorliegt;
        bei $\text{konst}/0$ gilt die Regel nicht.
  \item Ableiten von Zähler und Nenner getrennt — nicht die
        Quotientenregel verwenden.
  \item Die Regel ggf.\ mehrfach anwenden, aber nach jedem Schritt die Form
        neu prüfen.
  \item Alleine das Verschwinden von Zähler und Nenner am Punkt reicht; die
        Stetigkeit der Ableitungen ist nicht erforderlich.
\end{itemize}
\end{achtung}
